\documentclass[../../../include/open-logic-section]{subfiles}

\begin{document}

\olfileid{sth}{choice}{tarskiscott}
\olsection{તાર્સ્કી--સ્કોટની યુક્તિ}

\olref[cardinals][cardsasords]{defcardinalasordinal}માં આપણે
ગણાંકોને ક્રમવાચક સંખ્યાઓ તરીકે વ્યાખ્યાયિત કર્યા.
એમ કરવા આપણે સુક્રમિતતાની સ્વયંસિદ્ધિ માની હતી.
એવું કરવાનું એકમાત્ર કારણ એ હતું કે આ
``ક્ષેત્રની પ્રચલિત રીત'' છે.

તેથી સુક્રમિતતા વિશેના તત્ત્વચિંતનાત્મક પ્રશ્નો
ચર્ચીએ તે પહેલાં એ સ્પષ્ટ કરવું જરૂરી છે કે આપણે
પ્રચલિત રીતથી \emph{અલગ} જઈ \emph{સુક્રમિતતા
માન્યા વિના} પણ ગણાંકોનો સિદ્ધાંત વિકસાવી
\emph{શકીએ} છીએ. આપણે
\olref[sfr][siz][]{chap}માં આવેલી
$\cardeq{A}{B}$, $\cardle{A}{B}$ અને $\cardless{A}{B}$ની
વ્યાખ્યાઓ હજી પણ વાપરી શકીએ છીએ. ફક્ત
\emph{ગણાંક}નો નવો ખ્યાલ જોઈએ.

એક ભોળો વિચાર એ છે કે $A$ની ગણાંકતા આ રીતે
વ્યાખ્યાયિત કરીએ:
\[
	\Setabs{x}{\cardeq{A}{x}}.
\]
તેની સરખામણી \olref[cardinals][hp]{sec}ના
અંતમાં રૂપરેખા આપેલી ફ્રેગેની $\# x Fx$ની
વ્યાખ્યા સાથે કરી શકો. ત્યાં સૂચવેલાં કારણોને
લીધે આ વ્યાખ્યા નિષ્ફળ જાય છે. એક ઘટકવાળો
કોઈ પણ ગણ $\{\emptyset\}$ સાથે સમકદ છે.
પણ પદાનુક્રમના દરેક અનુગામી તબક્કે નવા
એક-ઘટકી ગણ રચાય છે (પાછલા તબક્કાનો
એક-ઘટકી ગણ જુઓ). તેથી
$\Setabs{x}{\cardeq{A}{x}}$ અસ્તિત્વમાં નથી,
કારણ કે તેનો સ્તરાંક હોઈ શકે નહીં.

આ સમસ્યા ટાળવા તાર્સ્કી અને સ્કોટની એક
યુક્તિ વાપરીએ:\footnote{યાદ રાખો: વિપરીત
સ્પષ્ટપણે ન કહ્યું હોય તો બધાં સૂત્રોમાં
પરિમાણો હોઈ શકે છે.}

\begin{defn}[તાર્સ્કી--સ્કોટ]
કોઈ પણ સૂત્ર $\phi(x)$ માટે $[ x : \phi(x)] $
એ શક્ય હોય એવા સૌથી નાના સ્તરાંકવાળા,
$\phi(x)$ને સંતોષતા બધા $x$નો ગણ છે
(અને $\phi$ને સંતોષતું કંઈ ન હોય તો $\emptyset$).
\end{defn}

આ વ્યાખ્યા યોગ્ય છે તે તપાસીએ. $\ZF$માં
કામ કરતાં,
\olref[spine][foundation]{zfentailsregularity}
દરેક $x$ માટે $\setrank{x}$નું અસ્તિત્વ
સુનિશ્ચિત કરે છે. હવે $\phi$ને સંતોષતી
કોઈ વસ્તુઓ હોય, તો $\alpha$ને એવો સૌથી
નાનો સ્તરાંક લઈએ કે
$(\exists x\subseteq V_\alpha)\phi(x)$,
એટલે કે $(\forall \beta \in \alpha)(\forall x
\subseteq V_\beta)\lnot \phi(x)$.
પછી પૃથક્કરણથી $[x : \phi(x)]$ને
$\Setabs{x \in V_{\alpha+1}}{\phi(x)}$
તરીકે વ્યાખ્યાયિત કરી શકીએ.

તાર્સ્કી--સ્કોટની યુક્તિને વાજબી ઠરાવ્યા
પછી તેના વડે ગણાંકતાનો એક ખ્યાલ
વ્યાખ્યાયિત કરી શકીએ:

\begin{defn}
$A$ની \textsc{ts}-ગણાંકતા
$\text{tsc}(A) = [x : \cardeq{A}{x}]$ છે.
\end{defn}

\textsc{ts}-ગણાંકની વ્યાખ્યા સુક્રમિતતા
વાપરતી નથી. પરંતુ એ સ્વયંસિદ્ધિ માન્યા
વિના પણ આપણે બતાવી શકીએ કે
\emph{\textsc{ts}-ગણાંકો},
\olref[cardinals][cardsasords]{defcardinalasordinal}માં
વ્યાખ્યાયિત \emph{ગણાંકો}ની જેમ ઘણુંખરું
વર્તે છે. ઉદાહરણ તરીકે,
\olref[cardinals][cardsasords]{lem:CardinalsBehaveRight}
અને \olref[card-arithmetic][opps]{lem:SizePowerset2Exp}ને
\textsc{ts}-ગણાંકોની ભાષામાં ફરી કહીએ, તો
સુક્રમિતતા માન્યા વિના $\ZF$માં તેમના
પુરાવા સફળ થાય છે.

આ જ પ્રસંગે એ પણ નોંધવું જોઈએ કે
તાર્સ્કી--સ્કોટની યુક્તિથી ક્રમવાચક
સંખ્યાઓનો સિદ્ધાંત પણ વિકસાવી શકીએ.
$\tuple{A, <}$ સુક્રમ હોય, ત્યારે
$\text{tso}(A, <) = [\tuple{X,
R} : \ordeq{\tuple{A, <}}{\tuple{X, R}}]$ લો.
ગણાંકો અને ક્રમવાચક સંખ્યાઓ અંગેના આ
અભિગમ માટે \citet[chs.~9--12]{Potter2004} જુઓ.

\end{document}
